\section{Mission and communities}\label{sec:paul-mission}

Paul's mission is often drawn as arrows radiating from one heroic traveler.
His letters show something more social. He entered cities through existing
roads and relationships, worked and lodged with others, proclaimed Christ,
formed assemblies in households, moved on before every problem was solved,
and then governed at a distance through messengers, letters, return visits,
delegates, and prayer. The communities were not trophies. They were agents
whose gifts, conflicts, endurance, and material support carried the mission
beyond Paul's presence.

\subsection{Antioch, Barnabas, and an expanding field}

Galatians places Paul in Syria and Cilicia after his first Jerusalem visit and
later associates him closely with Barnabas. Acts supplies the fuller account:
Barnabas brings Saul from Tarsus to Antioch; both teach there; the church sends
them with relief for Judea and then, under the Spirit, commissions Barnabas
and Saul with John Mark. They travel through Cyprus and southern Anatolia,
proclaiming first in synagogues and then among Gentiles. On Cyprus Acts begins
to call Saul by the name Paul. In Pisidian Antioch, Iconium, Lystra, and Derbe
it narrates
mixed response, signs, violent opposition, and the appointment of elders.

The journey belongs to Acts' first missionary circuit. Paul's letters do not
recount it as a unit. Galatians does confirm mission in Galatia, a severe
illness at the first proclamation, the Galatians' welcome, and a subsequent
crisis over circumcision. Whether ``Galatia'' names the southern Roman
provincial territory visited in Acts 13--14 or ethnic Galatia farther north
remains debated and affects the letter's date.

Acts follows the mission with the Jerusalem deliberation over Gentile
believers and then a sharp separation: Barnabas wants to take Mark; Paul
refuses because Mark had withdrawn earlier; Barnabas sails to Cyprus with
Mark, while Paul leaves with Silas. Galatians confirms tension at Antioch and
Barnabas's involvement in Peter's withdrawal but does not narrate the later
split. Neither source permits Barnabas to vanish into the role of a failed
assistant. He had advocated for Paul, shared danger and mission, and continued
a ministry Acts no longer follows.

\subsection{Macedonia and Achaia}

Acts' next circuit brings Silas, Timothy, and eventually the narrator's ``we''
to Macedonia. It names Lydia, a dealer in purple goods, as the first baptized
household at Philippi; narrates imprisonment, hymns, earthquake, and the
jailer's household; and then follows conflict in Thessalonica and Berea.
Paul's First Letter to the Thessalonians supplies the earliest surviving
inside view. The community turned from idols, endured pressure, worried about
members who had died, and became bound to Paul, Silvanus, and Timothy through
labor, affection, and hope. Paul compares his care both to a nursing mother
and to an exhorting father.

At Athens, Acts presents the Areopagus speech as a model of address to pagan
hearers through altar, creation, poetry, judgment, and resurrection. Paul
himself later recalls arriving at Corinth in weakness, fear, and trembling,
determined to know Christ crucified. The emphases are different, not mutually
exclusive. Athens is a crafted Lukan speech; Corinth is Paul's rhetorical
memory within a community dispute.

Corinth became the best-documented Pauline church. Paul worked with Aquila and
Priscilla according to Acts, stayed through Gallio's proconsulship, and later
answered reports and questions about faction, sex, marriage, lawsuits,
idol-food, worship, gifts, Resurrection, and the Jerusalem collection. Visits
and letters multiplied. Second Corinthians remembers a painful visit, a
severe or tearful letter, reconciliation through Titus, continuing rivalry,
and uncertainty over the collection. The correspondence reveals pastoral
authority exercised amid mutual hurt, not effortless control.

\subsection{Ephesus and Asia}

Acts places Paul in Ephesus for more than two years and makes the city a hub
from which the word spread through the province of Asia. It reports teaching,
healings, failed magical imitation, burning of books, and a theater riot
arising from the silver trade associated with Artemis. First Corinthians,
written from Ephesus, independently discloses a wide field: many adversaries,
an ``open door,'' daily danger, co-workers moving between cities, and a plan to
remain until Pentecost. Paul says he fought with beasts at Ephesus, but whether
this is literal arena exposure or metaphor for violent opponents is disputed.

Communities in the Lycus valley---Colossae, Laodicea, and Hierapolis---belong
to the Pauline mission's wider network even where Paul may not have founded or
visited them. Colossians and Philemon name Epaphras and a household circle;
the disputed authorship of Colossians means its biographical use must remain
qualified. Paul was not the sole traveler. Priscilla and Aquila, Apollos,
Timothy, Titus, Epaphroditus, Phoebe, and many unnamed bearers carried
teaching, money, letters, and news.

\subsection{The collection for Jerusalem}

The multi-province collection for poor believers in Jerusalem occupied years
of Paul's work. Galatians remembers the pillars' request that he remember the
poor. First Corinthians establishes a weekly saving practice; Second
Corinthians negotiates readiness, generosity, equality, accountability, and a
delegation; Romans asks prayer that the service be accepted. The gift was
material relief and a sign that Gentile communities shared in Israel's
spiritual blessings. It was also risky: money creates suspicion, and Paul's
anxiety about reception shows that unity could not be announced into existence.

Acts narrates Paul's arrival with representatives from the mission and then
his arrest, but it does not describe a ceremonial reception of the collection.
Silence does not prove rejection. It leaves the outcome less visible than
Paul's investment in it.

\subsection{Travel, suffering, and signs}

Second Corinthians gives the severest first-person inventory: imprisonments,
beatings, lashings, stoning, shipwrecks, danger from rivers, robbers, cities,
wilderness, sea, and false brothers, as well as hunger, cold, and daily
pressure for the churches. Written before Acts' final shipwreck, it proves how
selective any narrative itinerary must be. Paul turns the catalogue against
boasting in strength. Apostolic credibility appears in endurance and weakness,
not invulnerability.

Acts reports exorcisms, healings, extraordinary signs, rescue from prison,
survival after stoning, and the raising of Eutychus. Paul refers more tersely
to signs and wonders and to the power of the Spirit. This publication neither
supplies naturalistic mechanisms nor transfers these reports to later relics.
In both corpora, signs belong to proclamation and divine action, while Paul's
body remains vulnerable.

\evidenceaxis{Paul built and sustained a translocal mission through communities
and co-workers from Syria to Achaia and Asia.}{The seven undisputed letters,
read beside Acts 11--21 and the Gallio anchor.}{Acts' three circuits are a
useful theological itinerary, not a complete travel diary; several locations,
visits, letter dates, and reported miracles have source-specific limits.}
